%======================================================================
\subsection{BW-SIM-04 --- Built, 1 Nov: Netgen in a worker process with a time limit}\label{sec:sim-mesh-worker-built}
%======================================================================

Netgen runs inside one C++ call and cannot be interrupted. On the T\_HUB\,=\,1.9 bladed impeller of the VL-P01 turbojet
(read from STEP, second order, MaxSize 2.0, MinSize 0.1, Moderate) it never returned: 1\,h 45\,min in ``Meshing\ldots''
until the process was killed. The volume mesh now runs in a child process that the study can stop.

\paragraph{How it works.} The module is \file{bruceworx\_sim\_meshworker.py}. The shape is written as a BREP and the
settings (the Netgen properties of the study's mesh object, including \code{MeshSizeFile} of the local sizes) as JSON into a
temporary directory; a child \code{FreeCADCmd} of the same application (start-up 0.3\,s) meshes them and writes
\file{mesh.unv}, which keeps second order (Tet10 reads back as 10 nodes). The parent polls the child every 50\,ms and kills
its whole process group on timeout or cancel. Netgen holds the interpreter lock, so the child cannot notice a dead
parent itself: a small shell watchdog kills the group when the parent has gone (the app was quit during a mesh). The
temporary directory is removed in every case.

\paragraph{Time limit.} One try may take \code{MeshTimeLimit} seconds (a study property, 0 = automatic); automatic is
$\max(120\,\mathrm{s},\ 0.0007\,\mathrm{s}\times\hat N)$ capped at 3600\,s, $\hat N$ the node estimate of the pre-check. The
rule is 20 times what Netgen needs on the T\_HUB\,=\,1.6 impeller (17\,s for 514\,000 nodes). The tries of one meshing share
a total of four limits; a try gets what is left of it.

\paragraph{When the worker is used.} When the study sets a \code{MeshTimeLimit}, or the estimate is 20\,000 nodes or more.
Smaller models mesh in this process as before (the worker costs about a second), so a plate with a hole or a cantilever are
unchanged; the fallbacks of commit \code{6d41badd} (repaired copy, first order) remain the path there.

\paragraph{The tries (worker path).}
\begin{enumerate}
  \item Netgen at the asked order. The shape reaches the worker as a BREP, which is the ``repaired copy'' of the in-process
        fallback (a BREP round trip), so the separate repaired-copy try would repeat try 1 and is not run.
  \item When second order was asked: first order in a worker; the existing step upgrades the linear tetrahedra to Tet10.
  \item Gmsh, when a \code{gmsh} binary is found (FEM preference, next to the application, \code{PATH}). \textbf{Neither
        the build nor the packaged \file{Vector Labs.app} has one}, so today this try does not exist; the code is written
        (\code{gmsh -3 -order N -clmax MaxSize}, .unv back) but untested.
\end{enumerate}
Each try reports in the report view, the run log and the Progress tab (\code{bruceworx\_sim\_events} message at its end,
a progress text every 5\,s with the elapsed time), and \code{stats['notes']} names the failed tries and the one that worked.
If every try fails the error lists each with its time.

\paragraph{Cancel.} \code{bruceworx\_sim\_meshworker.cancel()} kills the running worker from any thread; the study ends with
``Meshing cancelled'' and no further try. The button in the Mesh node and the study panel is \textbf{not wired}: the panes
file belonged to another session at the time. It is one call.

\paragraph{Acceptance: the T\_HUB\,=\,1.9 impeller.} Read from the STEP (5.5\,MB, one solid), limit 300\,s per try:
try 1 (second order) no mesh after 300\,s, worker killed (300.1\,s); try 2 (first order) no mesh after 300\,s, killed (300.0\,s);
no Gmsh; the error lists both; 600\,s in all, then the program is free again (before: 1\,h 45\,min and a kill by hand).
Nothing is left behind. The same at MaxSize 3.0 (a coarser mesh) also gave 300.1\,s and 300.1\,s without a mesh, so
this solid is not a matter of element count: Netgen stalls in its surface or volume step on this geometry. Whether a
longer limit would finish was not tried (the first run was 1\,h 45\,min without an answer).

\paragraph{Tests} (\code{TestBruceWorxSimRotor.TestMeshWorker}, headless, 7): the worker mesh equals the in-process mesh on a
60\,$\times$\,10\,$\times$\,10 cantilever (nodes, volumes, volume, 4 and 10 nodes per tetrahedron); a worker that sleeps past
its limit is killed and the first-order try then meshes (upgraded to Tet10); every try timing out raises an error naming both;
cancel kills the worker; the shell watchdog kills a child whose parent has gone; the time-limit rule; a small model stays
in-process. Every test checks that no \code{bw-meshjob-*} directory and no worker process remains. Measured on the cantilever:
worker mesh 0.3\,s of start-up plus the mesh. \textbf{Open questions.} (1) The hang is Netgen's on this solid at these
settings; the real fix is geometry (T\_HUB, root fillets) or a coarser size, not a faster mesher: is a Gmsh binary worth
adding to the bundle? (2) Local-size model tests (\code{TestBruceWorxSimModel}, 4 failures) belong to the SIM-07 P2 work in the tree and were failing without this change. (3) The Cancel button position (Mesh node, study panel, or the run ribbon)? (4) Should the automatic
limit be shown in the Mesh node so it can be set there?

\paragraph{Your answers (1 Nov).} Add Gmsh as a second mesher, bundled in the app and tried automatically when Netgen
fails or times out, also selectable per study. While a mesh or solve runs, the Compute button becomes Cancel (and each
running Progress row gets a stop button). The Mesh node shows the time limit, automatic or your own.
